\capitulo{2}{Objetivos del proyecto}

El objetivo principal de este TFG es desarrollar un sistema que registre, analice y visualice la interacción del usuario durante el proceso de modelado 3D usando Blender.

Se trata de crear un instrumento que permita registrar y organizar los datos generados durante el modelado en forma automática.

Además, nuestro objetivo es analizar la forma en que el usuario modela, cuales son las herramientas más utilizadas y cómo evoluciona el modelo a lo largo de cada sesión.

Nuestro sistema intenta recabar información relevante sobre el proceso de modelado sin alterar la experiencia del usuario. El propósito es obtener datos que ayuden a entender cómo se elaboran los modelos 3D, las tareas principales que se realizan y qué ocurre con la interacción dentro del espacio tridimensional.

No solo se recoge información sino que además se analiza y representa visualmente. Por ello, el proyecto incluye funciones que faciliten la vista y el análisis de la información recopilada desde distintas perspectivas: observación de las acciones secuenciales, variaciones en la escena, acontecimientos significativos, métricas visualizadas como gráficos y vistas 3D.

Con estos recursos, el proyecto se adapta al análisis de interacción en entornos tridimensionales y las estadísticas educativas en la aplicación de herramientas digitales.

\section{Objetivos específicos}

Para conseguir todo esto, nos hemos marcado estos objetivos más concretos:

\begin{itemize}

\item Diseñar un sistema para recoger datos que no moleste nada al usuario y que funcione sin que se dé cuenta mientras modela en 3D.

\item Crear el primer \textit{addon} para Blender que pueda guardar información sobre el tiempo, el espacio y la estructura de la escena y los objetos que se están modelando.

\item Que el sistema detecte automáticamente cosas importantes que pasan al modelar, como cuando se duplican objetos, se unen vértices, se cambia el modo de edición o se tocan las coordenadas UV (que dicen dónde van las texturas en el modelo 3D).

\item Asegurarnos de que los datos guardados no se pierdan entre una sesión de trabajo y otra, usando un formato bien organizado para poder analizarlos después.

\item Incluir una forma clara para que el usuario dé su permiso, así la recogida de datos será ética y transparente.

\item Comprobar que el sistema funciona bien haciendo pruebas y dejando que lo usen otras personas.

\item Recoger datos de verdad usando una versión ya estable del sistema en sesiones de modelado.

\item Desarrollar un segundo \textit{addon} que sirva para analizar y mostrar los datos que se han recogido, mirando el tiempo, el espacio y las estadísticas.

\end{itemize}

Con estos objetivos, podemos organizar el proyecto en varias etapas, desde cuando implementamos el sistema que guarda los datos hasta que analizamos y mostramos la información que obtenemos al usarlo en la vida real.